Suppose $\pi=\dfrac{a}{b}$ for positive integers $a,b$. For a positive integer $n$, define
\[
f(x)=\dfrac{x^n(a-bx)^n}{n!},\qquad
F(x)=f(x)-f''(x)+f^{(4)}(x)-\cdots+(-1)^n f^{(2n)}(x).
\]
Every derivative of $f$ at $0$ is an integer. Indeed, the coefficient of $x^{n+j}$ is $(-1)^j\binom nj a^{n-j}\dfrac{b^j}{n!}$, and multiplication by $(n+j)!$ makes it integral; other derivatives vanish. Also $f(\pi-x)=f(x)$, so derivatives at $\pi$ are integers as well. Hence $F(0),F(\pi)$ are integers.

The alternating sum gives $F''+F=f$, and therefore
\[
\int_0^\pi f(x)\sin x\,dx
=[F'(x)\sin x-F(x)\cos x]_0^\pi=F(\pi)+F(0),
\]
an integer. The integrand is positive on $(0,\pi)$. On the other hand,
\[
0<\int_0^\pi f(x)\sin x\,dx
\le\dfrac{\pi}{n!}\left(\dfrac{a^2}{4b}\right)^n,
\]
because $x(a-bx)\le \dfrac{a^2}{(4b)}$. For sufficiently large $n$, the right side is less than $1$ since factorials outgrow any fixed exponential (successive ratios eventually become less than $\dfrac{1}{2}$). This produces an integer strictly between $0$ and $1$, a contradiction. Thus $\pi$ is irrational.
